

%\setcounter{page}{\value{page}}
\chapter{Introducere}
% \label{cap:Introducere}
%\setcounter{page}{7}
%\addtocounter{page}{6}

Limbajele de programare compilate, precum \textbf{C} și \textbf{C++}, sunt de cele mai multe ori
alegerea corectă pentru dezvoltarea de sisteme și aplicații pentru care performanța în timpul
rulării este un aspect de importanță critică, însă alegerea lor prezintă inconveniente în ceea ce
privește procesul de build. Cerințele ridicate de memorie și timpii îndelungați de compilare cresc
proporțional cu dimensiunea proiectului aflat în dezvoltare.

Timpii îndelungați de compilare sunt de natură să incetinească drastic procesul de dezvoltare, ceea
ce determină întârzieri ce afectează livrarea produsului către clienți și pierderi costisitoare de
timp pentru evoluția proiectului.

Evident, investiția în hardware mai performant pentru a reduce timpul de compilare poate diminua
magnitudinea problemei, însă se estimează că în următorii zece ani, puterea de procesare se va satura
(legea lui Moor nu va mai fi valabilă începând din 2036 \cite{powell2008quantum}), iar accentul se 
va muta asupra optimizării resurselor de rețea. Progresele în domeniul comunicațiilor vor permite
transferuri rapide și eficiente de date între nodurile sistemelor distribuite, susținând astfel 
implementarea și extinderea proceselor distribuite de compilare. Această abordare permite
fragmentarea și procesarea paralelă a sarcinilor de build pe mai multe noduri din rețea, fie ele stații locale, servere dedicate sau resurse din cloud. Drept rezultat, proiectele software vor putea beneficia de timpi reduși de compilare și de o utilizare mai eficientă a resurselor disponibile, contribuind la creșterea productivității și la reducerea blocajelor din ciclurile de dezvoltare.

Compilarea folosind mai multe sisteme de calcul care comunică prin rețea este susținută și de faptul
că, odată cu răspândirea utilizării generale a calculatoarelor și a resurselor cloud, timpii de
neutilizare ai sistemelor (Idle Times), au crescut considerabil. În instituții educaționale, calculatoarele au o durată de neutilizare medie de 28\% (adică 2672 de ore din total de 9550 de 
ore de funcționare) \cite{kondo2009cost}. Așadar, exploatarea resurselor de calcul ale sistemelor
neutilizate reprezintă o eficientizare binevenită pentru reducerea duratei de build.

% \section{Fundamentare teoretică}
Pentru a putea înțelege de ce durata procesului de compilare pentru \textbf{C/C++} (dar mai ales pentru C++) este semnificativ mai mare față de alte limbaje de programare, găsesc esențial să 
explicăm modul de desfășurare al acestui proces.

Transformarea codului sursă într-un fișier executabil este compusă din trei etape majore. Primul
pas îl reprezintă expandarea directivelor de preprocesare, moment în care dependințele dintre
fișierele sursă ale proiectului sunt determinate. Rezultatul etapei de preprocesare îl reprezintă
o colecție de unități de translație (Translation Units), fiecare fișier sursă având asociat un
astfel de \textbf{translation unit}. Urmează etapa de compilare, în urma căreia unitățile de
translație sunt transformate în fișiere ce conțin cod mașină, numite fișiere obiect (Object Files).
În ultima etapa, numită link-editare, toate fișierele obiect rezultate în urma compilării sunt 
conectate pentru a genera fișierul exectuabil al proiectului \cite{stroustrup2013c++}.

% \section{Preprocesarea}

Etapa de preprocesare este \cite{stroustrup2013c++} responsabilă de pregătirea fișierelor sursă
pentru a fi furnizate ca intrare pentru compilator. În \textbf{C/C++}, fiecare simbol poate fi definit o singură dată, însă declararea simbolului poate să apară de mai multe ori, întrucât fiecare 
simbol utilizat într-o unitate de translație, trebuie să aibă o declarație locală. Această chestiune
este tratată de mecanismul de includere a fișierelor, folosind directiva de preprocesare
\textbf{\#include}. Figura \ref{fig:preprocessing} exemplifică mecanismul de preprocesare al
fișierelor sursă.

\begin{figure}[H]
    \centering
    \includegraphics[width=0.8\linewidth, height=0.60\linewidth]{preprocessing.jpg}
    \caption{Preprocesarea fișierelor sursă \cite{stroustrup2013c++}}
    \label{fig:preprocessing}
\end{figure}

% \subsection{Compilarea unităților de translație}

Unitățile de translație sunt furnizate ca date de intrare pentru compilator, care detectează
erorile de sintaxă și produce fișierele obiect, care deși conțin cod mașină, nu pot fi încă executate. Figura exemplifică etapa de compilare.

\begin{figure}[H]
    \centering
    \includegraphics[width=0.8\linewidth, height=0.50\linewidth]{compilation.jpg}
    \caption{Compilarea unităților de translație \cite{stroustrup2013c++}}
    \label{fig:compilation}
\end{figure}

% \subsubsection{Link-editarea}

Etapa de linking \cite{stroustrup2013c++} conectează fișierele obiect pentru ca declarațiile locale
să fie disponibile în  celelalte unități de translație. Ieșirea linker-ului reprezintă un fișier 
executabil, care este rezultatul final al procesului de build. 

\begin{figure}[H]
    \centering
    \includegraphics[width=0.8\linewidth, height=0.55\linewidth]{linking.jpg}
    \caption{Generearea fișierului executabil \cite{stroustrup2013c++}}
    \label{fig:linking}
\end{figure}

% \subsection{Cauze pentru creșterea duratei de build}

Proiectele complexe, în care numărul de dependințe este mare, suferă des de recompilarea unui număr
mare de fișiere, ceea ce crește semnificativ durata de compilare. Motivul pentru care compilarea
excesivă este greu de evitat îl reprezintă faptul că, dacă un fișier header care este inclus de
un număr mare de fișiere sursă este modificat, atunci se declanșează recompilarea tuturor fișierelor
sursă care depind de acel header. Când numărul de fișiere header incluse este mare, există șansa
să se creeze cicluri de dependință, o situație în care mai multe fișiere se includ reciproc
\cite{martin2000design}.

Totodată, mecanismul de \textbf{templating} din \textbf{C++}, permite definirea de clase și
funcții generice, care pot fi instanțiate folosind tipuri de date diverse. Întrucât \textbf{C++}
este un limbaj de programare cu sistem static de tipuri (statically typed), de fiecare data
când un template este invocat pentru un nou tip de argument, codul pentru implementarea specifică
acelui tip de date este generat în mod automat de către compilator \cite{herity1998c++}. În cazul în care template-ul este instanțiat pentru foarte multe tipuri de date diferite, compiler-ul generează un volum semnificativ de cod repetitiv, ceea ce poate încetini procesul de build 
\cite{rotithor1999measurement}. 

Pentru a atrage atenția asupra scenariilor în care timpul de build devine o problemă, voi invoca
creșterea exponențială a volumului de cod sursă care intră în componența unor proiecte open-source
cunoscute, precum \textbf{Linux Kernel} și \textbf{Android}\cite{burnette2009hello}. Dimensiunea
codebase-ului pentru acestea se ridică la ordinul zecilor de milioane de linii de cod. Creșterea
volumului de cod de-a lungul timpului este vizibilă în Figura \ref{fig:locs}.

\begin{figure}[H]
    \centering
    \includegraphics[width=0.65\linewidth, height=0.27\linewidth]{locs.jpg}
    \caption{Creșterea codebase-ului de-a lungul anilor pentru câteva proiecte open-source \cite{lim2014distributed}}
    \label{fig:locs}
\end{figure}

\newpage

Pentru sistemul de operare mobil \textbf{Android 4.2.2}, timpul de build total ajunge la 51 de
minute \cite{lim2014distributed}.

\begin{figure}[H]
    \centering
    \includegraphics[width=0.65\linewidth, height=0.32\linewidth]{android.jpg}
    \caption{Timp de build Android \cite{lim2014distributed}}
    \label{fig:android}
\end{figure}
